\chapter{Creating Repositories}\label{chap: Creating Repositories}

Creating a repository is simple.

\begin{itemize}
    \item Login to your account on GitHub.
    \item On the top menu-bar, towards the right side, Click the $+$ (plus) button.
    \item Click ``New repository''.
\end{itemize}

On the following screen, fill in details specific to the new repository.

\begin{itemize}
    \item Give the repository a name. You will be told if you have used it previously.
    \item Give a description too. Be meaningful. The description here will be written into a \path{README.md} file if you choose to have one.
    \item Make the repository visible as required.
    \item Add a README if required. The default is not to have one, but they are very handy, and, will be used to display details of the repository on its home page.
    \item The \path{.gitignore} file is a good one to have as it stops files you might not want to be uploaded for public visibility, from being uploaded. Depending on how your project will be generated of compiled, you can choose a suitable \path{.gitignore} from the drop-down list.
    \item Add a licence\footnote{Or License if you spell things the American way!}---there are may to choose from.
    \item Click ``Create repository''.
\end{itemize}

That's it. All you need to do now is:

\begin{itemize}
    \item Add and edit files, online, in the GitHub page for the repository.
    \item It is much easier, for anything more than quick changes, to clone the repository to your own PC and work on it there, pushing committed changes back as and when appropriate. See \sectref{sect: Cloning a repository} for details.
\end{itemize}

